%!TEX root = ../main.tex

% ========================================================================
%  ANEXOS - LISTADOS DE CODIGO
%  Este fichero se incluye despues de \appendix en main.tex.
%  Los \lstinputlisting apuntan a la carpeta Codigo/. Coloca tu codigo en
%  Memoria/Codigo/ (o ajusta la ruta) con esta estructura:
%     Codigo/Transmisor/   -> los .ino y .h del transmisor
%     Codigo/Receptor/     -> FlightInstruments_Receiver.ino
%     Codigo/PFD/          -> los .pde del Processing
%  Los estilos Arduino y Processing estan definidos en packages.tex.
% ========================================================================

\chapter{Código del firmware del transmisor}
\label{anexo:fw_tx}

Se incluye a continuación el código completo del firmware del nodo transmisor,
desarrollado en Arduino/C++ y descrito en el \cref{ch:firmware}. El sketch se
organiza en varias unidades de compilación: la configuración y las estructuras de
datos compartidas, la adquisición y calibración de cada sensor, la fusión de
actitud y el sistema de telemetría.

\section{Config.h}
\lstinputlisting[style=Arduino]{Codigo/Transmisor/Config.h}

\section{SensorData.h}
\lstinputlisting[style=Arduino]{Codigo/Transmisor/SensorData.h}

\section{FlightInstruments.ino}
\lstinputlisting[style=Arduino]{Codigo/Transmisor/FlightInstruments.ino}

\section{IMU.ino}
\lstinputlisting[style=Arduino]{Codigo/Transmisor/IMU.ino}

\section{Calibration.ino}
\lstinputlisting[style=Arduino]{Codigo/Transmisor/Calibration.ino}

\section{SensorFusion.ino}
\lstinputlisting[style=Arduino]{Codigo/Transmisor/SensorFusion.ino}

\section{Barometer.ino}
\lstinputlisting[style=Arduino]{Codigo/Transmisor/Barometer.ino}

\section{GPS.ino}
\lstinputlisting[style=Arduino]{Codigo/Transmisor/GPS.ino}

\section{ToF.ino}
\lstinputlisting[style=Arduino]{Codigo/Transmisor/ToF.ino}

\section{Radar.ino}
\lstinputlisting[style=Arduino]{Codigo/Transmisor/Radar.ino}

\section{Telemetry.ino}
\lstinputlisting[style=Arduino]{Codigo/Transmisor/Telemetry.ino}


\chapter{Código del firmware del receptor}
\label{anexo:fw_rx}

Se incluye el código del nodo receptor, que decodifica los paquetes ARINC 429
recibidos por radio y los entrega al PC por el puerto serie, según lo descrito en
el \cref{ch:firmware}.

\section{FlightInstruments\_Receiver.ino}
\lstinputlisting[style=Arduino]{Codigo/Receptor/FlightInstruments_Receiver.ino}


\chapter{Código del PFD en Processing}
\label{anexo:pfd}

Se incluye el código completo del Primary Flight Display, desarrollado en
Processing y descrito en el \cref{ch:software}. Cada instrumento de la disposición
Basic-T se implementa en su propio fichero.

\section{FlightDisplay.pde}
\lstinputlisting[style=Processing]{Codigo/PFD/FlightDisplay.pde}

\section{AttitudeIndicator.pde}
\lstinputlisting[style=Processing]{Codigo/PFD/AttitudeIndicator.pde}

\section{HeadingIndicator.pde}
\lstinputlisting[style=Processing]{Codigo/PFD/HeadingIndicator.pde}

\section{AltimeterMSL.pde}
\lstinputlisting[style=Processing]{Codigo/PFD/AltimeterMSL.pde}

\section{GroundSpeedIndicator.pde}
\lstinputlisting[style=Processing]{Codigo/PFD/GroundSpeedIndicator.pde}

\section{VSI.pde}
\lstinputlisting[style=Processing]{Codigo/PFD/VSI.pde}

\section{ProximityIndicator.pde}
\lstinputlisting[style=Processing]{Codigo/PFD/ProximityIndicator.pde}

\section{GPSMap.pde}
\lstinputlisting[style=Processing]{Codigo/PFD/GPSMap.pde}